% Copyright (c) 2026 Intel Corporation.  All rights reserved.  See the file COPYRIGHT for more information.
% SPDX-License-Identifier: Apache-2.0

\documentclass[11pt]{article}

% Page layout
\usepackage[letterpaper, margin=1in]{geometry}
\setlength{\headheight}{14pt}
\usepackage{parskip}
\setlength{\parindent}{0pt}

% Typography
\usepackage{fancyvrb}
\usepackage{booktabs}
\usepackage{amsmath}
\usepackage{amssymb}

% Code listings
\usepackage{listings}
\lstset{
    basicstyle=\ttfamily\small,
    columns=fullflexible,
    keepspaces=true,
    showstringspaces=false,
    breaklines=true,
    frame=none,
    xleftmargin=2em,
}

\lstdefinelanguage{scheme}{
    morekeywords={define,let,let*,if,cond,lambda,begin,set!,and,or,not,
                  for-each,map,cons,car,cdr,null?,eq?,equal?,list,
                  BDD.New,BDD.And,BDD.Or,BDD.Not,BDD.True,BDD.False,
                  BDD.MakeTrue,BDD.MakeFalse,BDD.Equal,BDD.Equivalent,
                  BDD.Size},
    sensitive=true,
    morecomment=[l]{;},
    morestring=[b]",
}

% Hyperlinks
\usepackage[colorlinks=true, linkcolor=blue, urlcolor=blue, citecolor=blue]{hyperref}

% Headers and footers
\usepackage{fancyhdr}
\pagestyle{fancy}
\fancyhf{}
\fancyhead[L]{\textit{BDD-Based Symbolic Deadlock Checking}}
\fancyhead[R]{\thepage}
\renewcommand{\headrulewidth}{0.4pt}

\title{\textbf{BDD-Based Symbolic Deadlock Checking \\
for the Fulcrum CSP Toolchain} \\[0.5em]
\large Implementation Report}
\author{Mika Nystr\"om \and Claude (Anthropic)}
\date{March 2026}

\begin{document}

\maketitle

\tableofcontents
\newpage

%======================================================================
\section{Introduction}
%======================================================================

The Fulcrum CSP toolchain includes an explicit-state deadlock checker
(\texttt{check-deadlock!} in \texttt{product.scm}) that constructs the
product LTS of a system of communicating processes and explores it via
breadth-first search.  While correct, this approach enumerates every
reachable product state individually, and the state space grows
exponentially in the number of process instances.  For the probe-based
dining philosophers benchmark at $N=6$, the explicit checker explores
approximately 133,000~states in roughly 15~minutes.

This report describes the design and implementation of a complementary
\emph{symbolic} deadlock checker that represents state sets as Binary
Decision Diagrams (BDDs).  BDDs encode Boolean functions as directed
acyclic graphs with shared sub-structure, often compressing
exponentially large state sets into compact polynomial-size
representations.  The symbolic checker computes reachability via a
fixed-point iteration over BDD-encoded state sets, then checks for
deadlocked states---reachable states with no successor in the
transition relation.

The implementation adds a single new Scheme file
(\texttt{symbolic.scm}, $\sim$230 lines) to the CSP toolchain.  It
reuses the per-process LTS extraction pipeline and the system topology
parser, replacing only the product-state exploration engine.

%======================================================================
\section{Background}
%======================================================================

%----------------------------------------------------------------------
\subsection{Fulcrum CSP}
%----------------------------------------------------------------------

Fulcrum CSP~\cite{FulcrumCSP} is a programming language for
asynchronous circuit design, descended from Hoare's Communicating
Sequential Processes~\cite{Hoare85} via Martin's Communicating
Hardware Processes (CHP)~\cite{Martin90}.  Developed at Fulcrum Microsystems in the
early 2000s and subsequently used at Intel's Neuromorphic Computing
Lab, the language has seen over two decades of successful use in
hardware description, validation, and silicon compilation for dozens
of microchips, including products at the forefront of Ethernet
switching technology.

A Fulcrum CSP system consists of a static graph of independent
processes that do not share internal state but communicate exclusively
via the \emph{rendezvous}---a combined synchronization and
communication primitive.  When a sender reaches a send statement
\texttt{X!v} and the corresponding receiver reaches \texttt{Y?w} on
the connected channel, the distributed assignment $w \leftarrow v$ is
executed and both processes continue.  If either side arrives first, it
blocks until the other is ready.  This is the \emph{slack-zero}
semantics; channels may also have slack~$n > 0$, allowing the sender
to proceed up to $n$~times ahead of the receiver (FIFO
buffering)~\cite{Martin81}.

Two language features are particularly relevant to verification:

\textbf{Guarded commands.}  Following Dijkstra~\cite{Dijkstra75} and
Hoare, Fulcrum CSP uses guarded selections and repetitions for
branching and looping.  The selection
\texttt{[\,e -> S0\;[]\;f -> S1\,]} blocks until at least one guard
is true, then nondeterministically chooses among the true guards.  The
repetition \texttt{*[\,e -> S0\;[]\;f -> S1\,]} iterates until no
guard is true, then terminates.

\textbf{The probe.}  The probe operator \texttt{\#A}, introduced by
Martin, is a Boolean value that is true exactly when an action on
port~\texttt{A} can proceed without suspension.  Probes enable
reactive, nondeterministic process designs---a process can test which
of its neighbors is ready before committing to a communication---and
are essential for modeling hardware arbiters and multiplexers.  In the
LTS extraction pipeline, probes are abstracted as $\tau$~transitions
(internal nondeterministic choice), which is sound for deadlock
analysis.

Process creation and channel interconnection are handled outside the
CSP language itself, by a system description (\texttt{.sys}) file that
specifies process types, instances, channels, and port bindings.
This separation matches the hardware distinction between fabrication
(static topology) and operation (dynamic behavior).

%----------------------------------------------------------------------
\subsection{Explicit-State Deadlock Checking}
%----------------------------------------------------------------------

The existing checker (\texttt{product.scm}) works as follows:

\begin{enumerate}
\item Parse the \texttt{.sys} file to obtain the system topology:
      process types, instances, port bindings, and channels.
\item Extract a per-process-type LTS by compiling each CSP source file
      and traversing its blocking-point graph (\texttt{lts.scm}).
\item Rename port names to channel names according to the instance bindings.
\item Explore the product state space via BFS.  At each product state
      (a tuple of per-instance local states), compute successors:
      \begin{itemize}
        \item \emph{Tau interleaving}: any instance with a $\tau$ transition
              advances alone; all others remain.
        \item \emph{Channel synchronization}: if a sender instance can
              \texttt{send} on channel~$c$ and the corresponding receiver
              can \texttt{recv} on~$c$, both advance simultaneously.
      \end{itemize}
\item A product state with no successors is a deadlock.
\end{enumerate}

The state space is $\prod_{i} |S_i|$ where $S_i$ is the local state
set of instance~$i$.  For the $N$-philosopher dining system with
probe-based forks (8~local states per fork, 4 per philosopher), the
product space is $8^N \cdot 4^N = 32^N$, growing as $32^N$.  At
$N=8$ this is $32^8 \approx 10^{12}$---completely intractable for
explicit enumeration.

%----------------------------------------------------------------------
\subsection{Symbolic Model Checking with BDDs}
%----------------------------------------------------------------------

Symbolic model checking~\cite{McMillan93,Bryant86,ClarkeGrumberg99} represents sets of
states as Boolean functions encoded by BDDs.  A state is a bit-vector;
a set of states is the characteristic function
$f : \{0,1\}^n \to \{0,1\}$ that is~1 exactly for states in the set.
The transition relation $T(V, V')$ is a Boolean function over current-state
variables~$V$ and next-state variables~$V'$.

Key BDD operations used:
\begin{itemize}
\item $\mathit{And}(f, g)$, $\mathit{Or}(f, g)$, $\mathit{Not}(f)$:
      set intersection, union, complement.
\item $\exists x.\, f = \mathit{Or}(f|_{x=1},\; f|_{x=0})$:
      existential quantification (projection).
\item $f[v'/v]$: variable substitution (rename next-state to current-state).
\item $\mathit{Equal}(f, g)$: structural equality test (constant time
      with hash-consing).
\end{itemize}

Reachability is computed as a least fixed point:
\[
  R_0 = \mathit{Init}, \qquad
  R_{k+1} = R_k \cup \mathit{Image}(R_k, T)
\]
where $\mathit{Image}(R, T) = (\exists V.\, R(V) \wedge T(V, V'))[V'/V]$.
The iteration converges when $R_{k+1} = R_k$.

Deadlocked states are those reachable but lacking any successor:
\[
  D = R \wedge \neg (\exists V'.\, T(V, V'))
\]
If $D \neq \mathit{false}$, the system has a deadlock.

%----------------------------------------------------------------------
\subsection{The Dining Philosophers Problem}
%----------------------------------------------------------------------

The dining philosophers problem, introduced by Dijkstra~\cite{Dijkstra71}
and later formulated in CSP by Hoare~\cite{Hoare85}, is a canonical
benchmark for deadlock analysis.  $N$~philosophers sit around a circular
table, each sharing a fork with each neighbor.  A philosopher must
acquire both adjacent forks to eat, then release them.  If all
philosophers simultaneously pick up their left fork, none can acquire
the right, and the system deadlocks.

Our benchmark uses two process types from the Fulcrum CSP dialect.
The \textbf{naive philosopher} acquires its left fork, then its right,
eats (modeled as a pair of send acknowledgments), and repeats:

\begin{lstlisting}[frame=single,xleftmargin=1em,xrightmargin=1em]
int x;
*[ LG?x ; RG?x ; LT!1 ; RT!1 ]
\end{lstlisting}

Here \texttt{LG} (left-get) and \texttt{RG} (right-get) are input
ports that receive the fork; \texttt{LT} (left-return) and \texttt{RT}
(right-return) are output ports that return it.  The sequential
ordering \texttt{LG?x\,;\,RG?x} makes every philosopher reach for its
left fork first, creating the classic circular-wait deadlock.

The \textbf{probe-based fork} uses CSP channel probes for
nondeterministic selection between its two neighbors:

\begin{lstlisting}[frame=single,xleftmargin=1em,xrightmargin=1em]
int x;
*[[ #LG -> LG!1 ; LA?x
  : #RG -> RG!1 ; RA?x
 ]]
\end{lstlisting}

The guarded repetition \texttt{*[[\,\ldots\,]]} selects nondeterministically
among guards whose channel probes (\texttt{\#LG}, \texttt{\#RG}) are
true---i.e., the corresponding neighbor is ready to receive.  The fork
sends a token on the selected side (\texttt{LG!1} or \texttt{RG!1})
and waits for its return (\texttt{LA?x} or \texttt{RA?x}).  Despite
the nondeterminism, the system still deadlocks: when all philosophers
simultaneously request their left fork, every fork's \texttt{\#RG}
probe is false, so all forks commit to serving left, and the circular
wait forms.

For comparison, the \textbf{deterministic fork} always serves left
first, then right:

\begin{lstlisting}[frame=single,xleftmargin=1em,xrightmargin=1em]
int x;
*[ LG!1 ; LA?x ; RG!1 ; RA?x ]
\end{lstlisting}

The \textbf{reversed fork} (used in asymmetric configurations to break
circular wait) serves right first:

\begin{lstlisting}[frame=single,xleftmargin=1em,xrightmargin=1em]
int x;
*[ RG!1 ; RA?x ; LG!1 ; LA?x ]
\end{lstlisting}

The system topology is specified in a \texttt{.sys} file.  For $N=3$
with probe-based forks:

\begin{lstlisting}[frame=single,xleftmargin=1em,xrightmargin=1em,
    basicstyle=\ttfamily\footnotesize]
system DiningThree;
  var f0lg : channel(32); var f0la : channel(32);
  var f0rg : channel(32); var f0ra : channel(32);
  var f1lg : channel(32); var f1la : channel(32);
  var f1rg : channel(32); var f1ra : channel(32);
  var f2lg : channel(32); var f2la : channel(32);
  var f2rg : channel(32); var f2ra : channel(32);

  process Fork = "phil_fork.csp"
    port out LG : channel(32)  port in  LA : channel(32)
    port out RG : channel(32)  port in  RA : channel(32);

  process Phil = "phil_naive.csp"
    port in  LG : channel(32)  port out LT : channel(32)
    port in  RG : channel(32)  port out RT : channel(32);

begin
  var fork0 : Fork(LG=>f0lg, LA=>f0la, RG=>f0rg, RA=>f0ra);
  var fork1 : Fork(LG=>f1lg, LA=>f1la, RG=>f1rg, RA=>f1ra);
  var fork2 : Fork(LG=>f2lg, LA=>f2la, RG=>f2rg, RA=>f2ra);
  var phil0 : Phil(LG=>f0lg, LT=>f0la, RG=>f1rg, RT=>f1ra);
  var phil1 : Phil(LG=>f1lg, LT=>f1la, RG=>f2rg, RT=>f2ra);
  var phil2 : Phil(LG=>f2lg, LT=>f2la, RG=>f0rg, RT=>f0ra);
end.
\end{lstlisting}

Each philosopher's left fork is the fork with the same index; its
right fork is the next fork (modulo~$N$).  The port bindings connect
each philosopher's \texttt{LG}/\texttt{LT} to its left fork's
\texttt{LG}/\texttt{LA}, and \texttt{RG}/\texttt{RT} to the right
fork's \texttt{RG}/\texttt{RA}.  The pattern scales to arbitrary~$N$
by adding more fork/philosopher instance pairs and channels.

The \textbf{asymmetric} configuration, which is deadlock-free, replaces
one fork with the reversed variant (\texttt{phil\_fork\_rev.csp}),
breaking the circular wait.

%======================================================================
\section{Implementation}
%======================================================================

The symbolic checker is implemented in \texttt{csp/src/symbolic.scm}
and loaded from \texttt{cspc.scm} after \texttt{product.scm}.  It
uses the BDD library available via \texttt{import("bdd")} in the
\texttt{m3makefile}, with Scheme bindings generated by
\texttt{SchemeStubs("BDD")}.

%----------------------------------------------------------------------
\subsection{BDD Interface}
%----------------------------------------------------------------------

The BDD library provides the following Scheme-accessible operations:

\begin{center}
\begin{tabular}{ll}
\toprule
\textbf{Operation} & \textbf{Description} \\
\midrule
\texttt{(BDD.New \textit{name})} & Create a fresh BDD variable \\
\texttt{(BDD.True)}, \texttt{(BDD.False)} & Boolean constants \\
\texttt{(BDD.And \textit{a b})}, \texttt{(BDD.Or \textit{a b})} & Conjunction, disjunction \\
\texttt{(BDD.Not \textit{a})} & Negation \\
\texttt{(BDD.Equivalent \textit{a b})} & Bi-implication ($a \Leftrightarrow b$) \\
\texttt{(BDD.MakeTrue \textit{f v})} & Cofactor $f|_{v=1}$ \\
\texttt{(BDD.MakeFalse \textit{f v})} & Cofactor $f|_{v=0}$ \\
\texttt{(BDD.Equal \textit{a b})} & Structural equality (Boolean result) \\
\texttt{(BDD.Size \textit{a})} & Node count \\
\bottomrule
\end{tabular}
\end{center}

%----------------------------------------------------------------------
\subsection{State Encoding}
%----------------------------------------------------------------------

Each process instance~$i$ with $K_i$ local states is encoded using
$b_i = \lceil \log_2 K_i \rceil$ BDD variable pairs:

\begin{itemize}
\item \textbf{Current-state variables}: $v^c_{i,0}, v^c_{i,1}, \ldots, v^c_{i,b_i-1}$
\item \textbf{Next-state variables}: $v^n_{i,0}, v^n_{i,1}, \ldots, v^n_{i,b_i-1}$
\end{itemize}

Variables are created with names like \texttt{"i0\_c0"},
\texttt{"i0\_n0"} (instance~0, current bit~0 / next bit~0).

States are mapped to integers $0, 1, \ldots, K_i - 1$.  The encoding
of ``instance~$i$ is in state~$s$'' is the conjunction of variable
polarities matching the binary representation of~$s$:
\[
  \mathit{encode}(i, s) = \bigwedge_{k=0}^{b_i-1}
    \begin{cases}
      v^c_{i,k} & \text{if bit $k$ of $s$ is 1} \\
      \neg v^c_{i,k} & \text{otherwise}
    \end{cases}
\]

%----------------------------------------------------------------------
\subsection{Transition Relation}
%----------------------------------------------------------------------

The transition relation $T(V, V')$ is the disjunction of all individual
transition BDDs:

\textbf{Tau transition} of instance~$i$ from state~$s$ to state~$s'$:
\[
  \mathit{encode\_curr}(i, s) \wedge \mathit{encode\_next}(i, s')
  \wedge \bigwedge_{j \neq i} \mathit{frame}(j)
\]

\textbf{Channel synchronization} on channel~$c$ with sender~$i_s$
going $s \to s'$ and receiver~$i_r$ going $r \to r'$:
\[
  \mathit{encode\_curr}(i_s, s) \wedge \mathit{encode\_next}(i_s, s')
  \wedge \mathit{encode\_curr}(i_r, r) \wedge \mathit{encode\_next}(i_r, r')
  \wedge \bigwedge_{j \neq i_s, i_r} \mathit{frame}(j)
\]

The \textbf{frame condition} for instance~$j$ constrains its next
state to equal its current state:
\[
  \mathit{frame}(j) = \bigwedge_{k=0}^{b_j-1} (v^c_{j,k} \Leftrightarrow v^n_{j,k})
\]

%----------------------------------------------------------------------
\subsection{Image Computation}
%----------------------------------------------------------------------

The image of a state set $R$ under $T$ is computed in three steps:

\begin{enumerate}
\item \textbf{Conjoin}: $C = R(V) \wedge T(V, V')$
\item \textbf{Project}: $P = \exists V.\, C$ \quad (quantify out current-state variables)
\item \textbf{Rename}: $I = P[V' \mapsto V]$ \quad (shift next-state to current-state)
\end{enumerate}

Existential quantification is implemented as:
\[
  \exists x.\, f = \mathit{Or}(\mathit{MakeTrue}(f, x),\; \mathit{MakeFalse}(f, x))
\]
applied sequentially for each variable in~$V$.

Variable renaming substitutes each next-state variable $v'$ with the
corresponding current-state variable~$v$:
\[
  f[v'/v] = \mathit{Or}(\mathit{And}(v, \mathit{MakeTrue}(f, v')),\;
                         \mathit{And}(\neg v, \mathit{MakeFalse}(f, v')))
\]

%----------------------------------------------------------------------
\subsection{Helper Functions}
%----------------------------------------------------------------------

\begin{lstlisting}[language=scheme]
(define (bdd-exists bdd var)
  (BDD.Or (BDD.MakeTrue bdd var) (BDD.MakeFalse bdd var)))

(define (bdd-exists-list bdd vars)
  (if (null? vars) bdd
      (bdd-exists-list (bdd-exists bdd (car vars)) (cdr vars))))

(define (bdd-compose bdd old new)
  (BDD.Or (BDD.And new (BDD.MakeTrue bdd old))
           (BDD.And (BDD.Not new) (BDD.MakeFalse bdd old))))

(define (bdd-rename bdd old-vars new-vars)
  (if (null? old-vars) bdd
      (bdd-rename (bdd-compose bdd (car old-vars) (car new-vars))
                  (cdr old-vars) (cdr new-vars))))
\end{lstlisting}

%----------------------------------------------------------------------
\subsection{Main Algorithm}
%----------------------------------------------------------------------

The entry point \texttt{check-deadlock-symbolic!} performs:

\begin{enumerate}
\item Parse the \texttt{.sys} file and validate the system topology
      (reusing \texttt{parse-sys-file} and \texttt{validate-system!}
      from \texttt{cspbuild.scm}).
\item Extract per-process LTSs and rename ports to channels (reusing
      \texttt{extract-system-lts!} and \texttt{rename-lts} from
      \texttt{product.scm}).
\item Build BDD encodings for each instance (state-to-integer mapping,
      current/next variable pairs).
\item Build the channel map (sender/receiver index per channel).
\item Construct the monolithic transition relation $T(V, V')$.
\item Precompute the has-successor set:
      $H = \exists V'.\, T(V, V')$.
\item Run fixed-point reachability from the initial state.
\item Check: $D = R \wedge \neg H$.  If $D \neq \mathit{false}$,
      report deadlock.
\end{enumerate}

The function returns \texttt{\#t} (deadlock-free) or \texttt{\#f}
(deadlock exists).  Unlike the explicit checker, it does not currently
produce a counterexample trace, since extracting a witness path from
the BDD encoding requires additional machinery (e.g., backward
unrolling or on-the-fly trace recording).

%======================================================================
\section{Experimental Results}
%======================================================================

All experiments were run on an Apple M3 (ARM64\_DARWIN) with the
CM3 compiler and the BDD library from the \texttt{caltech-other/bdd}
package.  Timing is wall-clock time measured with \texttt{Time.Now}
from the Modula-3 runtime; it includes LTS extraction (parsing and
compiling the CSP source files), which is shared overhead for both
checkers.

%----------------------------------------------------------------------
\subsection{Small Systems: Correctness Validation}
%----------------------------------------------------------------------

We first validated that the symbolic checker agrees with the explicit
checker on all existing test systems:

\begin{center}
\begin{tabular}{lcccc}
\toprule
\textbf{System} & \textbf{Instances} & \textbf{Explicit} & \textbf{Symbolic} & \textbf{Agree?} \\
\midrule
\texttt{prodcons}         & 2  & deadlock-free (4 states)  & deadlock-free & \checkmark \\
\texttt{pipeline}         & 3  & deadlock-free (10 states) & deadlock-free & \checkmark \\
\texttt{dining\_det} (N=2)     & 4  & DEADLOCK (4 states)  & DEADLOCK & \checkmark \\
\texttt{dining\_asym} (N=2)    & 4  & deadlock-free (8 states)  & deadlock-free & \checkmark \\
\texttt{dining\_det} (N=3)     & 6  & DEADLOCK (8 states)  & DEADLOCK & \checkmark \\
\texttt{dining\_asym} (N=3)    & 6  & deadlock-free (33 states) & deadlock-free & \checkmark \\
\texttt{dining\_naive} (N=2)   & 4  & DEADLOCK (32 states) & DEADLOCK & \checkmark \\
\bottomrule
\end{tabular}
\end{center}

For these small systems both checkers take comparable wall-clock time
(1.3--4.2\,s), dominated by the shared LTS extraction overhead.

%----------------------------------------------------------------------
\subsection{Scaling: Probe-Based Dining Philosophers}
%----------------------------------------------------------------------

The probe-based dining philosophers system provides a challenging
scalability benchmark.  Each fork process uses channel probes
(\texttt{*[\,\#LG -> \ldots\;[\!]\;\#RG -> \ldots\,]}) for
nondeterministic selection, yielding 8~local states per fork and
4~per philosopher.  The total product state space is~$32^N$.

\begin{center}
\begin{tabular}{rrrrrrrr}
\toprule
$N$ & \textbf{Inst.} & \textbf{Bits} & $|\mathbf{S}|$ &
  \textbf{$|T|$ nodes} & \textbf{Peak $|R|$} & \textbf{Iter.} & \textbf{Time} \\
\midrule
2  &  4 & 10 & $10^3$      &   195 &    --- &  --- &   3.6\,s  \\
3  &  6 & 15 & $3\!\times\!10^4$ &   397 &  4,793 &  38 &  17\,s   \\
6  & 12 & 30 & $10^9$      & 1,363 &  4,793 &  38 &  17\,s   \\
8  & 16 & 40 & $10^{12}$   & 2,307 & 50,550 &  52 &  12\,min \\
\bottomrule
\end{tabular}
\end{center}

\begin{itemize}
\item \textbf{$|T|$ nodes}: size of the transition relation BDD.
      Grows linearly in $N$ (approximately $190 \cdot N$).
\item \textbf{Peak $|R|$}: maximum BDD node count for the reachable
      set during the fixed-point iteration.  This is the primary
      driver of memory and time.
\item \textbf{Iter.}: number of fixed-point iterations to convergence.
      Corresponds to the diameter of the reachable state graph.
\item \textbf{$|S|$}: total product state space ($32^N$).
      The explicit checker must enumerate all reachable states
      individually; at $N=6$ it explores 133K~states in $\sim$15\,min.
      At $N=8$ ($\sim\!10^{12}$ states) it is completely intractable.
\end{itemize}

The symbolic checker found DEADLOCK in all cases (as expected for the
naive philosopher protocol with symmetric forks).

%----------------------------------------------------------------------
\subsection{BDD Size Evolution}
%----------------------------------------------------------------------

The BDD size of the reachable set during the $N=8$ fixed-point
computation shows a characteristic pattern:

\begin{center}
\begin{tabular}{rr@{\qquad}rr@{\qquad}rr}
\toprule
\textbf{Iter.} & \textbf{BDD size} &
\textbf{Iter.} & \textbf{BDD size} &
\textbf{Iter.} & \textbf{BDD size} \\
\midrule
 0 &    102 & 18 & 18,653 & 36 & 43,531 \\
 5 &    541 & 20 & 27,410 & 40 & 40,546 \\
10 &  2,701 & 25 & 45,252 & 45 & 39,234 \\
15 & 10,985 & 29 & \textbf{50,550} & 50 & 39,010 \\
16 & 13,868 & 30 & 50,283 & 52 & 39,033 \\
\bottomrule
\end{tabular}
\end{center}

The BDD grows during the expanding wavefront phase (iterations 0--29),
peaks at $\sim$50K nodes, then \emph{shrinks} as the reachable set
fills in and regularities emerge that the BDD can exploit via
sub-graph sharing.  This is a well-known phenomenon in symbolic model
checking: the final fixed-point BDD is often smaller than intermediate
BDDs because the complete reachable set has more structure than partial
wavefronts.

%======================================================================
\section{Memory Analysis}
\label{sec:memory}
%======================================================================

The high memory consumption observed during the $N=8$ run is expected
and stems from several factors inherent to BDD-based symbolic model
checking.

%----------------------------------------------------------------------
\subsection{Sources of Memory Usage}
%----------------------------------------------------------------------

\textbf{1. The BDD unique table.}
BDD libraries maintain a hash table (the \emph{unique table}) that
ensures each Boolean function has exactly one canonical BDD node.
Every call to \texttt{BDD.And}, \texttt{BDD.Or}, etc.\ looks up or
creates entries in this table.  Nodes are never duplicated, but the
table itself can grow large.  At $N=8$ the reachable set peaks at
50K~nodes, but the unique table also retains all intermediate results
from prior iterations unless explicitly garbage-collected.

\textbf{2. Intermediate BDDs in image computation.}
Each image computation step involves:
\begin{enumerate}
\item $C = \mathit{And}(R, T)$: this conjunction can be much larger
      than either operand.  With $|R| \approx 50\text{K}$ and
      $|T| \approx 2.3\text{K}$, the conjunction $C$ can temporarily
      reach hundreds of thousands of nodes.
\item Sequential existential quantification: $\exists v_0.\,(\exists v_1.\,
      (\cdots(\exists v_{39}.\, C)\cdots))$.  Each of the 40
      quantification steps produces a new intermediate BDD.  While
      quantification generally \emph{reduces} the BDD size (by
      removing a variable), the intermediate results can temporarily
      blow up before the variable is fully eliminated.
\item Sequential variable renaming: 40 substitution steps, each
      creating a new intermediate BDD.
\end{enumerate}

At $N=8$ with 40~BDD variables, each fixed-point iteration performs
$\sim$80 intermediate BDD operations (40 quantifications + 40 renames),
times 52~iterations $= \sim$4,200 intermediate BDDs.  The unique table
accumulates nodes from all of these.

\textbf{3. No automatic garbage collection.}
The BDD library from \texttt{caltech-other/bdd} uses a simple
reference-based scheme.  Unlike BDD packages with built-in garbage
collection (e.g., CUDD's periodic dead-node reclamation), this library
retains all nodes created during the computation.  Dead nodes---those
unreachable from any live BDD root---are not reclaimed automatically.
Over 52~iterations, the accumulated dead nodes can significantly
exceed the live node count.

%----------------------------------------------------------------------
\subsection{Estimated Memory Footprint}
%----------------------------------------------------------------------

A rough estimate: each BDD node requires $\sim$20--32~bytes (variable
index, low/high child pointers, reference count, hash link).  If the
unique table accumulates $\sim$500K--1M nodes over the full
computation (including dead nodes from intermediate results), the BDD
storage alone would consume 15--30\,MB.  However, the Modula-3
runtime's traced heap and the Scheme interpreter's own allocations
(cons cells for the image computation loop, lambda closures, etc.) add
substantial overhead, and the total process memory can reach several
hundred megabytes.

%----------------------------------------------------------------------
\subsection{Mitigation Strategies}
%----------------------------------------------------------------------

Several well-known techniques can reduce memory consumption in future
work:

\begin{enumerate}
\item \textbf{BDD garbage collection.}  Periodically invoke a
      mark-and-sweep pass on the unique table, retaining only nodes
      reachable from live roots ($R$, $T$, $H$, and the variable
      nodes).  CUDD~\cite{Somenzi98} does this automatically.

\item \textbf{Early quantification.}  Instead of forming the full
      conjunction $R \wedge T$ and then quantifying out all
      current-state variables, interleave conjunction with
      quantification: quantify out variables as soon as they are no
      longer needed.  This keeps intermediate BDDs smaller.
      The technique is known as \emph{partitioned transition
      relations} or \emph{conjunctive partitioning}~\cite{BurchClarke91}.

\item \textbf{Partitioned transition relations.}  Instead of building
      a single monolithic $T$, maintain $T$ as a conjunction of
      per-instance transition BDDs.  During image computation, conjoin
      and quantify one partition at a time.  This avoids ever
      materializing the full $T \wedge R$ product.

\item \textbf{Variable reordering.}  Dynamic BDD variable reordering
      (sifting~\cite{Rudell93}) can dramatically reduce BDD sizes by
      finding a variable order that minimizes sharing.  The current
      implementation uses a fixed interleaved order
      ($v^c_{0,0}, v^n_{0,0}, v^c_{0,1}, v^n_{0,1}, \ldots$),
      which is a reasonable default but not necessarily optimal.

\item \textbf{Frontier simplification.}  Replace the reachable set $R$
      with a simplified ``frontier'' set $F_k = R_k \setminus R_{k-1}$
      (only newly discovered states), compute the image of $F_k$, and
      accumulate into $R$.  This reduces the size of the BDD entering
      the image computation, since $F_k$ is typically much smaller
      than~$R_k$.
\end{enumerate}

%======================================================================
\section{Integration}
%======================================================================

%----------------------------------------------------------------------
\subsection{Files Modified}
%----------------------------------------------------------------------

\begin{center}
\begin{tabular}{llp{7.5cm}}
\toprule
\textbf{File} & \textbf{Action} & \textbf{Description} \\
\midrule
\texttt{csp/src/symbolic.scm} & Created & BDD-based symbolic deadlock checker ($\sim$230 lines) \\
\texttt{csp/src/cspc.scm} & 1 line added & \texttt{(load "symbolic.scm")} after \texttt{product.scm} \\
\texttt{csp/src/m3makefile} & 1 line added & \texttt{SchemeStubs ("Time")} for benchmarking \\
\texttt{csp/src/tests/run\_deadlock\_tests.sh} & Batch 4 added & 6 symbolic checker tests (4 deadlock-free, 2 deadlock) \\
\bottomrule
\end{tabular}
\end{center}

%----------------------------------------------------------------------
\subsection{Test Results}
%----------------------------------------------------------------------

All 15 tests pass (9~existing explicit-state tests plus
6~new symbolic tests):

\begin{Verbatim}[fontsize=\small]
=== Product LTS Deadlock Checking Acceptance Test Suite ===

--- Batch 1: deadlock-free systems ---
  PASS  deadlock-free:prodcons
  PASS  deadlock-free:pipeline
  PASS  deadlock-free:multi_inst
  PASS  deadlock-free:bidir

--- Batch 2: expected-deadlock systems ---
  PASS  deadlock:deadlock_simple (trace=3 steps)
  PASS  deadlock:deadlock_cycle (trace=3 steps)

--- Batch 3: dining philosophers ---
  PASS  deadlock:dining_det_2 (trace=3 steps)
  PASS  deadlock:dining_det_3 (trace=4 steps)
  PASS  deadlock:dining_probe_2 (trace=7 steps)

--- Batch 4: symbolic deadlock checking ---
  PASS  sym-deadlock-free:prodcons
  PASS  sym-deadlock-free:pipeline
  PASS  sym-deadlock-free:dining_asym_2
  PASS  sym-deadlock-free:dining_asym_3
  PASS  sym-deadlock:dining_det_2
  PASS  sym-deadlock:dining_det_3

=== Results: 15/15 passed ===
\end{Verbatim}

%======================================================================
\section{Limitations and Future Work}
%======================================================================

\begin{enumerate}
\item \textbf{No counterexample trace.}  The symbolic checker returns
      only a Boolean verdict.  Extracting a concrete deadlock trace
      from BDDs requires backward image computation from the deadlock
      states, selecting one concrete state per step.  This is
      straightforward but not yet implemented.

\item \textbf{Monolithic transition relation.}  The current
      implementation builds a single BDD for the entire transition
      relation.  Partitioned (conjunctive/disjunctive) representations
      would reduce peak memory during image computation.

\item \textbf{No garbage collection.}  As discussed in
      Section~\ref{sec:memory}, the BDD library does not reclaim dead
      nodes.  Adding periodic garbage collection or switching to a
      library with built-in reclamation (e.g., CUDD) would
      significantly reduce memory consumption.

\item \textbf{Sequential quantification and renaming.}  The current
      \texttt{bdd-exists-list} and \texttt{bdd-rename} functions
      process variables one at a time.  Vectorized operations
      (simultaneous quantification over a set of variables) would be
      more efficient, both in time and in intermediate BDD sizes.

\item \textbf{Livelock and refinement checking.}  The symbolic
      framework extends naturally to other verification properties:
      livelock detection (fair reachability under strong fairness) and
      trace/failures refinement checking.  These are planned for
      Phase~4 of the formal verification roadmap.
\end{enumerate}

%======================================================================
\section{Conclusions}
%======================================================================

The BDD-based symbolic deadlock checker successfully complements the
existing explicit-state checker.  On the probe-based dining
philosophers benchmark:

\begin{itemize}
\item At $N=6$ (12~instances, $10^9$ potential states), the symbolic
      checker finds deadlock in \textbf{17\,seconds} versus
      $\sim$15\,minutes for explicit enumeration of 133K~reachable states.

\item At $N=8$ (16~instances, $10^{12}$ potential states), the symbolic
      checker finds deadlock in \textbf{12\,minutes} with a peak BDD
      size of only 50K~nodes.  Explicit enumeration is completely
      infeasible at this scale.
\end{itemize}

The implementation is compact ($\sim$230 lines of Scheme), reuses the
existing LTS extraction and system parsing infrastructure, and
produces results that are validated against the explicit checker on all
small test cases.  The primary limitation is memory consumption due to
accumulated dead BDD nodes, which can be addressed by garbage
collection and partitioned transition relations in future work.

\begin{thebibliography}{99}

\bibitem{FulcrumCSP}
M.~Nystr\"om,
``The Fulcrum CSP Language,'' reference manual (draft),
Intel Corporation, 2025.
\url{https://github.com/mikanystrom/intel-async/blob/main/async-toolkit/m3utils/m3utils/csp/doc/csp/csp.tex}

\bibitem{Hoare85}
C.~A.~R. Hoare,
\emph{Communicating Sequential Processes}.
Prentice Hall, 1985.

\bibitem{Roscoe97}
A.~W. Roscoe,
\emph{The Theory and Practice of Concurrency}.
Prentice Hall, 1997.

\bibitem{Dijkstra71}
E.~W. Dijkstra,
``Hierarchical ordering of sequential processes,''
\emph{Acta Informatica}, vol.~1, no.~2, pp.~115--138, 1971.

\bibitem{Dijkstra75}
E.~W. Dijkstra,
``Guarded commands, nondeterminacy, and formal derivation of programs,''
\emph{Comm.\ ACM}, vol.~18, no.~8, pp.~453--457, 1975.

\bibitem{Martin81}
A.~J. Martin,
``An axiomatic definition of synchronization primitives,''
\emph{Acta Informatica}, vol.~16, pp.~219--235, 1981.

\bibitem{Martin90}
A.~J. Martin,
``Programming in VLSI: from communicating processes to delay-insensitive
circuits,''
in \emph{UT Year of Programming Institute on Concurrent Programming},
C.~A.~R. Hoare, Ed.\
Addison-Wesley, 1990, pp.~1--64.

\bibitem{Bryant86}
R.~E. Bryant,
``Graph-based algorithms for Boolean function manipulation,''
\emph{IEEE Trans.\ Computers}, vol.~C-35, no.~8, pp.~677--691, 1986.

\bibitem{McMillan93}
K.~L. McMillan,
\emph{Symbolic Model Checking}.
Kluwer Academic, 1993.

\bibitem{BurchClarke91}
J.~R. Burch, E.~M. Clarke, D.~E. Long,
``Symbolic model checking with partitioned transition relations,''
in \emph{Proc.\ VLSI}, 1991, pp.~49--58.

\bibitem{ClarkeGrumberg99}
E.~M. Clarke, O.~Grumberg, D.~A. Peled,
\emph{Model Checking}.
MIT Press, 1999.

\bibitem{Somenzi98}
F.~Somenzi,
``CUDD: CU Decision Diagram Package,''
University of Colorado at Boulder, 1998.

\bibitem{Rudell93}
R.~Rudell,
``Dynamic variable ordering for ordered binary decision diagrams,''
in \emph{Proc.\ ICCAD}, 1993, pp.~42--47.

\end{thebibliography}

\end{document}
